\documentclass[
10pt, % Main document font size
a4paper, % Paper type, use 'letterpaper' for US Letter paper
oneside, % One page layout (no page indentation)
%twoside, % Two page layout (page indentation for binding and different headers)
headinclude,footinclude, % Extra spacing for the header and footer
BCOR5mm, % Binding correction
]{scrartcl}
%%%%%%%%%%%%%%%%%%%%%%%%%%%%%%%%%%%%%%%%%
% Arsclassica Article
% Structure Specification File
%
% This file has been downloaded from:
% http://www.LaTeXTemplates.com
%
% Original author:
% Lorenzo Pantieri (http://www.lorenzopantieri.net) with extensive modifications by:
% Vel (vel@latextemplates.com)
%
% License:
% CC BY-NC-SA 3.0 (http://creativecommons.org/licenses/by-nc-sa/3.0/)
%
%%%%%%%%%%%%%%%%%%%%%%%%%%%%%%%%%%%%%%%%%

%----------------------------------------------------------------------------------------
%	REQUIRED PACKAGES
%----------------------------------------------------------------------------------------

\usepackage[
nochapters, % Turn off chapters since this is an article        
beramono, % Use the Bera Mono font for monospaced text (\texttt)
eulermath,% Use the Euler font for mathematics
pdfspacing, % Makes use of pdftex’ letter spacing capabilities via the microtype package
dottedtoc % Dotted lines leading to the page numbers in the table of contents
]{classicthesis} % The layout is based on the Classic Thesis style

\usepackage{arsclassica} % Modifies the Classic Thesis package
\usepackage{ulem}
\usepackage{graphicx}
\usepackage{float}
\usepackage{algorithm2e}
\usepackage[export]{adjustbox}
\usepackage[utf8]{inputenc}  
\usepackage[english,frenchb]{babel} 
\usepackage[T1]{fontenc} % Required for accented characters
\usepackage{multicol}
\usepackage{graphicx} % Required for including images
\graphicspath{{Figures/}} % Set the default folder for images

\usepackage{enumitem} % Required for manipulating the whitespace between and within lists

\usepackage{lipsum} % Used for inserting dummy 'Lorem ipsum' text into the template


\expandafter\def\csname ver@subfig.sty\endcsname{}


\usepackage{svg}
\usepackage{hyperref}

\usepackage{amsmath,amssymb,amsthm} % For including math equations, theorems, symbols, etc

\usepackage{varioref} % More descriptive referencing

%----------------------------------------------------------------------------------------
%	THEOREM STYLES
%---------------------------------------------------------------------------------------

\theoremstyle{definition} % Define theorem styles here based on the definition style (used for definitions and examples)
\newtheorem{definition}{Definition}

\theoremstyle{plain} % Define theorem styles here based on the plain style (used for theorems, lemmas, propositions)
\newtheorem*{remark}{Remark}
\newtheorem{theorem}{Theorème}
\newtheorem{lemma}[theorem]{Lemme}
\newtheorem{proposition}[theorem]{Proposition}

\theoremstyle{remark} % Define theorem styles here based on the remark style (used for remarks and notes)

%----------------------------------------------------------------------------------------
%	HYPERLINKS
%---------------------------------------------------------------------------------------

\hypersetup{
%draft, % Uncomment to remove all links (useful for printing in black and white)
colorlinks=true, breaklinks=true, bookmarks=true,bookmarksnumbered,
urlcolor=webbrown, linkcolor=RoyalBlue, citecolor=webgreen, % Link colors
pdftitle={}, % PDF title
pdfauthor={\textcopyright}, % PDF Author
pdfsubject={}, % PDF Subject
pdfkeywords={}, % PDF Keywords
pdfcreator={pdfLaTeX}, % PDF Creator
pdfproducer={LaTeX with hyperref and ClassicThesis} % PDF producer
}

\hyphenation{Fortran hy-phen-ation} 

\title{\normalfont{FTML Exercices 6 solutions}\\
} 
\author{\spacedlowsmallcaps{}
} 
\date{\today}
\begin{document}
\renewcommand{\sectionmark}[1]{\markright{\spacedlowsmallcaps{#1}}}
\lehead{\mbox{\llap{\small\thepage\kern1em\color{halfgray} \vline}\color{halfgray}\hspace{0.5em}\rightmark\hfil}} 
\pagestyle{scrheadings}
\maketitle 
\setcounter{tocdepth}{3} 

\tableofcontents

\section{\large\color{Blue}Jacobiennes}

\paragraph{f1}%
\label{par:f1}

Pour cet exemple le plus simple est d'utiliser les dérivées partielles. Si on note $J_{x, y, z}$ la jacobienne en $(x, y, z)$, on a:

\begin{equation}
J_{x, y, z} = 
\begin{pmatrix}
    -y\sin(xy) &  -x\sin (xy) & 0\\
    -z\sin(xz) &  0 & -x\sin(xz)\\
\end{pmatrix}
\end{equation}

\paragraph{f2}%
\label{par:f1}

Si on note $J_{a, b, c, d}$ la jacobienne en $(a, b, c, d)$, on a:

\begin{equation}
J_{a, b, c, d} = 
\begin{pmatrix}
    1 &  -1 & 0 & 0 \\
    1 &  -1 & -1 & -1 \\
    0 &  0 & 1 & 1 \\
    0 &  1 & 0 & -1 \\
\end{pmatrix}
\end{equation}


\paragraph{3}%
\label{par:f2}

Supposons par l'absurde que la norme $x\mapsto ||x||$ soit différentiable en $0$. Il existe donc $L$ linéaire continue telle que

\begin{equation}
    \label{eq:diff}
    ||h|| = L(h) + o(h)
\end{equation}

On peut réécrire $o(h)$ comme le produit $||h||a(h)$, avec $a$ une application tendant vers $0$ quand $h$ tend vers $0$. Soit $t$ un nombre réel. Avec l'homogénéité de la norme, on sait que $||th||=|t|||h||$. En appliquant \ref{eq:diff}, on obtient que
\begin{equation}
    \label{eq:homo}
    L(th) + ||th||a(th) = |t| (L(h) + ||h||a(h))
\end{equation}

Prenons $t>0$. Avec la linéarité de $L$, \ref{eq:homo} devient

\begin{equation}
    tL(h) + t||h||a(th) = t L(h) + t||h||a(h)
\end{equation}

Puis en simplifiant

\begin{equation}
    a(th) = a(h)
\end{equation}

En faisant tendre $t$ vers 0, on obtient que nécessairement $a(h)=0$ pour tout $h$. Ainsi, $||h||=L(h)$.

Par ailleurs on sait que si $h\neq 0$, alors $||h||>0$, et $||-h||>0$. Or, $||-h|| = L(h) = -L(h) = -||h||<0$, ce qui est absurde.


\end{document}
